% =====================================================================
\section{Conclusion}
\label{sec:conclusion}
% =====================================================================

Five findings show how the Euronext--TSO alliance affects Nord Pool.

\begin{itemize}
  \item \textbf{Ecosystem coordination} between the exchange and grid operators depends more on regulation than on ownership. Although two TSOs withdrew from the ownership structure in 2022, CACM and the shared coupling maintain their connection to the market and allow it to keep functioning.
  \item \textbf{Network effects:} the strongest effect, market coupling, is a regulated commons shared with competitors. What distinguishes Nord Pool now is the advantage provided by its system price, data and connection to futures.
  \item \textbf{Platform governance} gives the owners unequal influence: Euronext holds 66\% and chairs the board, while the TSOs occupy one of eight seats.
  \item \textbf{Cybersecurity} follows weakest-link provision at the interfaces, and much of the loss falls on actors who do not make investment decisions. The pricing that expands the network also adds weak entry points.
  \item \textbf{Value creation and value capture} are now separated. Trading fees represent roughly 0.15\% of the power price, with value capture increasingly coming from data and derivatives.
\end{itemize}

Our analysis identifies a central tension between using the price as a public signal and treating it as a private asset. The platform publishing area prices and calculating the Nordic reference price is controlled by a listed exchange group. State-owned TSOs supply the transmission capacities needed for area prices and bear the system risk, yet have limited influence over platform operations. This tension is visible in the introduction of the data paywall in 2022 \citep{ERR2022} and the launch of system-price futures in 2026 \citep{Euronext2026}. Members, TSOs and regulators create value together, but this value is increasingly sold through data and futures products within the Euronext group, while the division of revenue between Euronext and Nord Pool is not public.

The existence of this tension does not imply that the alliance destroys value. Capital and technology have come from Euronext, which has also taken over the Nordic power futures market from Nasdaq. Competition through the multi-NEMO arrangement may also help keep trading fees low, although we cannot establish this causally (Section~\ref{sec:t2}). The change is in how the public interest is protected: regulation now carries more of the responsibility previously held through ownership. Three developments will shape whether the alliance retains its legitimacy:
\begin{enumerate}
  \item whether regulators treat price data as public infrastructure;
  \item whether EPEX continues to provide credible competition;
  \item whether the TSOs retain sufficient influence to protect security of supply as their ownership role becomes less central, despite their essential contribution.
\end{enumerate}
Our essay examines these questions further through its broader economic and societal assessment.
